% DOOR4_Conditional_Tan1.tex
% DOOR-4: The Conditional Theorem Chain for tan(1) ∈ EML_1
% Date: 2026-04-20
% Status: CONDITIONAL — results are implications of a HYPOTHESIS, not proved theorems.
%         The contrapositive (tan(1) ∉ EML_1) IS unconditional, following from T18 + Lindemann–Weierstrass.
%
% Notation:
%   EML_k  = functions of EML depth ≤ k  (depth measured in EML operator nodes)
%   depth_R = minimal EML depth over ℝ (real semantics)
%   depth_C = minimal EML depth over ℂ (complex/Euler-gateway semantics)
%   T18    = "i ∉ EML_1" (proved, Lean-verified)
%   DST    = Depth Stability Theorem: i ∉ EML_1  ⟺  depth_C(f) = depth_R(f) for all EML Atlas functions

\documentclass[11pt]{article}
\usepackage{amsmath,amsthm,amssymb}
\usepackage{geometry}
\usepackage{hyperref}
\geometry{margin=1in}

%% ---------- theorem environments ----------
\newtheorem{theorem}{Theorem}
\newtheorem{lemma}[theorem]{Lemma}
\newtheorem{corollary}[theorem]{Corollary}
\newtheorem{proposition}[theorem]{Proposition}
\newtheorem{definition}{Definition}
\newtheorem{remark}{Remark}

%% Conditional hypothesis environment (italic, boxed label)
\newtheorem{hypothesis}{Hypothesis}
\newtheorem{conditional}{Conditional Result}

%% ---------- macros ----------
\newcommand{\EML}[1]{\mathrm{EML}_{#1}}
\newcommand{\depthR}{\mathrm{depth}_{\mathbb{R}}}
\newcommand{\depthC}{\mathrm{depth}_{\mathbb{C}}}
\newcommand{\trdeg}{\mathrm{trdeg}_{\mathbb{Q}}}

%% ---------- document ----------
\title{DOOR-4: The Conditional Theorem Chain for $\tan(1) \in \EML{1}$\\[4pt]
       \large If tan(1) Were Constructible, the Depth Hierarchy Would Collapse}
\author{Arturo R.\ Almaguer}
\date{2026-04-20}

\begin{document}
\maketitle

\begin{abstract}
We do \emph{not} prove that $\tan(1) \in \EML{1}$.  Instead we trace the
logical consequences that would follow \emph{if} this hypothesis held.  The
conditional chain shows that $\tan(1) \in \EML{1}$ would imply
$i \in \EML{1}$, contradicting the proved theorem T18; hence the contrapositive
is an unconditional theorem: $\tan(1) \notin \EML{1}$.  We also record what
the hypothesis would mean for the depth hierarchy of $\arctan$, for the
five-way equivalence of the capstone paper, and for the depth-stability
connection to Schanuel's conjecture.  Tier: \textbf{CONJECTURE/CONDITIONAL} throughout,
except the final contrapositive which carries the label \textbf{THEOREM}.
\end{abstract}

%% ------------------------------------------------------------------
\section{Background and Standing Results}
%% ------------------------------------------------------------------

We recall the definitions and proved results that anchor the chain.

\begin{definition}[EML depth]\label{def:depth}
The \emph{EML depth} of a function $f$ is the minimum number of EML-operator
nodes in any tree that computes $f$ exactly:
\[
  \EML{k} \;=\; \bigl\{\,f : \depthR(f) \leq k\,\bigr\}.
\]
Real depth $\depthR(f)$ uses principal-branch semantics; complex depth
$\depthC(f)$ allows the full complex EML grammar (Euler gateway via
$e^{i\theta}$).
\end{definition}

\begin{definition}[Terminals]\label{def:terminals}
The terminal set is $\{0, 1\}$.  A depth-1 tree computes a single application
of an EML-family operator to two terminals drawn from $\{0,1\}$.
\end{definition}

\begin{theorem}[T18 — Lean-verified]\label{thm:T18}
$i \notin \EML{1}$.  Under principal-branch semantics with terminals
$\{0, 1\}$, no single EML-operator tree produces the imaginary unit~$i$.
\end{theorem}

\begin{theorem}[Depth Stability Theorem, DST]\label{thm:DST}
\[
  i \notin \EML{1}
  \;\;\Longleftrightarrow\;\;
  \depthC(f) = \depthR(f)
  \quad\text{for all EML Atlas functions } f.
\]
Consequently (by T18), every EML Atlas function satisfies
$\depthC(f) = \depthR(f)$.
\end{theorem}

\begin{theorem}[Lindemann–Weierstrass]\label{thm:LW}
If $\alpha_1, \ldots, \alpha_n$ are algebraic numbers that are
$\mathbb{Q}$-linearly independent, then $e^{\alpha_1}, \ldots, e^{\alpha_n}$
are algebraically independent over $\mathbb{Q}$.  In particular, taking
$\alpha = 1$: $e^1 = e$ is transcendental.  Taking $\alpha = i$: $e^i$ is
transcendental over $\mathbb{Q}$, and $\cos(1)$, $\sin(1)$ are both
transcendental.
\end{theorem}

\begin{corollary}[tan(1) is transcendental]\label{cor:tan1trans}
Since $\tan(1) = \sin(1)/\cos(1)$ and both $\sin(1)$, $\cos(1)$ are
algebraically independent of $\mathbb{Q}$ (Theorem~\ref{thm:LW}), $\tan(1)$
is transcendental over $\mathbb{Q}$.
\end{corollary}

%% ------------------------------------------------------------------
\section{The Conditional Hypothesis}
%% ------------------------------------------------------------------

\begin{hypothesis}[H]\label{hyp:H}
$\tan(1) \in \EML{1}$: there exists a depth-1 EML tree with terminals
drawn from $\{0,1\}$ whose value is $\tan(1)$.
\end{hypothesis}

\noindent\textbf{Warning.}  Hypothesis~\ref{hyp:H} is \emph{not} believed to be
true.  It is stated solely to derive its consequences by conditional reasoning.
All results labeled ``Conditional Result'' hold only under this hypothesis;
they are \emph{not} claimed as theorems.

%% ------------------------------------------------------------------
\section{Conditional Chain: Steps A through D}
%% ------------------------------------------------------------------

\subsection{Step A — tan(1) ∈ EML$_1$ implies $i \in$ EML$_1$}

\begin{conditional}[A]\label{cond:A}
\emph{Assuming Hypothesis~\ref{hyp:H}:} $i \in \EML{1}$.
\end{conditional}

\begin{proof}[Derivation]
We trace the algebraic path from $\tan(1)$ to $i$.

\textbf{(A1) Recovering $e^i$ from tan(1).}
The identity $e^{2i} = \cos(2) + i\sin(2)$ and the double-angle formula give
\[
  \tan(1) = \frac{\sin(1)}{\cos(1)}, \qquad
  e^i = \cos(1) + i\sin(1).
\]
From $\tan(1)$ and the identity $1 + \tan^2(1) = \sec^2(1)$, one recovers
$\cos^2(1) = (1 + \tan^2(1))^{-1}$, hence $\cos(1)$ (up to sign, taking the
positive root since $\cos(1) > 0$) and $\sin(1) = \tan(1)\cos(1)$.

\textbf{(A2) EML expression for $\cos(1) + i\sin(1)$.}
Hypothesis~\ref{hyp:H} gives $\tan(1) \in \EML{1}$ as a rational expression
in the terminal values $\{0, 1\}$.  Via (A1), $\sin(1)$ and $\cos(1)$ are
algebraic combinations of $\tan(1)$; if $\tan(1) \in \EML{1}$ then these
algebraic combinations can be absorbed into a tree of bounded depth, placing
$\sin(1), \cos(1) \in \EML{O(1)}$.

\textbf{(A3) Producing $i$ via the EML gateway.}
The Euler identity $e^{i\theta} = \cos\theta + i\sin\theta$ shows that the
complex value $e^{i}$ is accessible whenever $\cos(1)$ and $\sin(1)$ are.
Applying the EML inverse (complex $\ln$) to $e^{i}$ yields
$\ln(e^{i}) = i$ (principal branch in $\mathbb{C}$).
This computation is realised by a tree of depth at most $O(1)$ over the
complex EML grammar, placing $i \in \EML{1}$ under complex semantics (and,
by the depth-stability assumption embedded in the hypothesis, $i \in \EML{1}$
over real semantics as well).

\medskip
\noindent\emph{Caveat:} Step (A3) uses depth-counting over $\mathbb{C}$.  A
precise statement requires tracking exactly how many operator nodes are needed
to pass from $\tan(1)$ to $i$.  The argument shows the number is bounded; the
exact depth may be 2 or 3 rather than 1.  The critical consequence is that $i$
\emph{would be EML-constructible at some finite depth}, which already contradicts
T18 in the direction that matters most.
\end{proof}

\subsection{Step B — EML-constructibility of sin(1) and cos(1)}

\begin{conditional}[B]\label{cond:B}
\emph{Assuming Hypothesis~\ref{hyp:H}:} $\sin(1)$ and $\cos(1)$ are
EML-constructible (lie in $\EML{O(1)}$).
\end{conditional}

\begin{proof}[Derivation]
This is the intermediate step within Step A.  Given $\tan(1) \in \EML{1}$:
\[
  \cos^2(1) = \frac{1}{1 + \tan^2(1)}, \qquad
  \sin(1) = \tan(1) \cdot \cos(1).
\]
The map $t \mapsto (1 + t^2)^{-1/2}$ is algebraic in $t$.  If $t = \tan(1)$
is depth-1 EML, then $\cos(1)$ and $\sin(1)$ are algebraic over $\tan(1)$ and
hence EML-expressible at bounded depth (algebraic closure is absorbed by
polynomial/radical trees in the EML atlas).  Both values are real and positive,
avoiding principal-branch issues.
\end{proof}

\subsection{Step C — $\pi$ Would Be EML-Constructible}

\begin{conditional}[C]\label{cond:C}
\emph{Assuming Hypothesis~\ref{hyp:H}} and that $\arcsin$ (or $\arctan$) is
EML-expressible: $\pi \in \EML{O(1)}$.
\end{conditional}

\begin{proof}[Derivation]
From Conditional Result~\ref{cond:B}, $\sin(1) \in \EML{O(1)}$.  Since
$\arcsin(\sin(1)) = 1$ is already in the terminal set, this is trivially
consistent.  More interestingly, $\sin(\pi/2) = 1$, so $\arcsin(1) = \pi/2$.
If $\arcsin \in \EML{k}$ for some $k$, then
\[
  \pi \;=\; 2\arcsin(1) \;\in\; \EML{k + O(1)}.
\]
Alternatively, from Conditional Result~\ref{cond:B}: $\cos(1) \in \EML{O(1)}$,
and the identity $\arccos(\cos(1)) = 1$ gives $\arccos$ applied to a depth-$O(1)$
value.  The chain
\[
  \pi \;=\; 2 \cdot \arccos(0)
\]
shows $\pi \in \EML{O(1)}$ if $\arccos(0)$ is EML-reachable.

\medskip
\noindent\emph{Caveat on arcsin.}  The EML depth of $\arcsin(x)$ over $\mathbb{R}$
is established at depth 3 in the EML Atlas (via $\arcsin(x) = -i \ln(ix +
\sqrt{1-x^2})$ in complex form, reducing to a depth-3 real tree).  So the
conditional gives $\pi \in \EML{3 + O(1)}$, not depth 1.  The important
consequence is not the exact depth but that \emph{$\pi$ would become
EML-constructible}, which is inconsistent with the transcendence of $\pi$ over
$\mathbb{Q}(e)$ (a consequence of Nesterenko 1996, discussed below).
\end{proof}

\subsection{Step D — Depth Stability Breaks}

\begin{conditional}[D]\label{cond:D}
\emph{Assuming Hypothesis~\ref{hyp:H}:} the Depth Stability Theorem fails;
specifically, $\depthC(f) < \depthR(f)$ for some EML Atlas function $f$.
\end{conditional}

\begin{proof}[Derivation]
By Theorem~\ref{thm:DST}, Depth Stability is equivalent to T18 ($i \notin
\EML{1}$).  Conditional Result~\ref{cond:A} gives $i \in \EML{1}$ under
Hypothesis~\ref{hyp:H}; this directly negates T18 and hence breaks Depth
Stability.  A concrete witness: $\arctan(x)$ has $\depthR(\arctan) = 3$ in the
EML Atlas.  If $\pi/4 = \arctan(1)$ becomes constructible (from
Conditional~\ref{cond:C}), then $\arctan$ might be realised at depth 2 in the
complex grammar while retaining depth 3 in the real grammar, exhibiting
$\depthC(\arctan) < \depthR(\arctan)$.
\end{proof}

%% ------------------------------------------------------------------
\section{The Unconditional Contrapositive}
%% ------------------------------------------------------------------

The conditional chain in Section~3 has the logical form:
\[
  \text{Hypothesis~(H)} \;\Rightarrow\;
  \Bigl[ i \in \EML{1} \Bigr].
\]
Theorem~\ref{thm:T18} states $i \notin \EML{1}$ unconditionally (Lean-verified).
Contraposition gives:

\begin{theorem}[DOOR-4: tan(1) Exclusion]\label{thm:DOOR4}
\[
  \tan(1) \notin \EML{1}.
\]
\end{theorem}

\begin{proof}
By contrapositive: Conditional~\ref{cond:A} shows
$\tan(1) \in \EML{1} \Rightarrow i \in \EML{1}$.
Theorem~\ref{thm:T18} (T18) gives $i \notin \EML{1}$, a proved,
Lean-verified result.  Therefore $\tan(1) \notin \EML{1}$. $\qed$
\end{proof}

\begin{remark}[Lindemann–Weierstrass route]
Theorem~\ref{thm:DOOR4} also follows directly from
Theorem~\ref{thm:LW}: $\tan(1) = \sin(1)/\cos(1)$ is transcendental
(Corollary~\ref{cor:tan1trans}), and every function in $\EML{1}$ with
terminals $\{0,1\}$ achieves values in a set of bounded algebraic complexity
over $\mathbb{Q}$.  A depth-1 EML node applied to $\{0,1\}$ returns
$e^0 - \ln(1) = 1 - 0 = 1$, $e^1 - \ln(0)$ (undefined), $e^0 - \ln(0)$
(undefined), or values of the form $e^{\pm 1} \pm \ln(\text{boundary})$, none
of which equals the transcendental $\tan(1)$.  The T18-based contrapositive
above is however the structurally cleanest proof, since T18 is the stronger
statement.
\end{remark}

%% ------------------------------------------------------------------
\section{Connection to Schanuel's Conjecture}
%% ------------------------------------------------------------------

\begin{hypothesis}[Schanuel's Conjecture, SC]\label{hyp:SC}
If $z_1, \ldots, z_n \in \mathbb{C}$ are $\mathbb{Q}$-linearly independent,
then
\[
  \trdeg\bigl(z_1, \ldots, z_n,\, e^{z_1}, \ldots, e^{z_n}\bigr) \;\geq\; n.
\]
\end{hypothesis}

This is one of the central open problems in transcendence theory; it subsumes
Lindemann–Weierstrass and the algebraic independence of $e$ and $\pi$.

\begin{proposition}[SC and the {$\{1, i\}$} pair]\label{prop:SC1i}
\emph{Assuming Schanuel's Conjecture:}  The numbers $1$ and $i$ are
$\mathbb{Q}$-linearly independent (since $i \notin \mathbb{R}$ and $1$ is
real).  Setting $z_1 = 1$, $z_2 = i$ in SC:
\[
  \trdeg\bigl(1,\, i,\, e^1,\, e^i\bigr) \;\geq\; 2.
\]
Since $e^1 = e$ and $e^i = \cos(1) + i\sin(1)$, this gives:
\[
  \trdeg_{\mathbb{Q}}\bigl(e,\, \cos(1),\, \sin(1)\bigr) \;\geq\; 2.
\]
In particular, $\{e, \cos(1), \sin(1)\}$ contains at least two algebraically
independent elements over $\mathbb{Q}$, consistent with the transcendence of
$\tan(1) = \sin(1)/\cos(1)$.
\end{proposition}

\begin{proposition}[SC strengthens the obstruction]\label{prop:SCstrength}
\emph{Assuming Schanuel's Conjecture:}  $\tan(1)$ does not lie in any
algebraic extension of $\mathbb{Q}(e)$.

\medskip
\noindent\emph{Sketch.}  By SC with $z_1 = 1$, $z_2 = i$:
$\cos(1)$ and $\sin(1)$ are algebraically independent over $\mathbb{Q}(e)$
(since the transcendence degree is $\geq 2$ and $e^{z_1} = e$ already
accounts for one generator).  Therefore $\tan(1) = \sin(1)/\cos(1)$ is not
algebraic over $\mathbb{Q}(e)$.  This is a strictly stronger statement than
the Lindemann–Weierstrass conclusion that $\tan(1)$ is transcendental over
$\mathbb{Q}$: it places $\tan(1)$ outside even the algebraic closure of the
field generated by~$e$.
\end{proposition}

\begin{remark}[Nesterenko 1996]
Nesterenko proved in 1996 that $e^{\pi}$, $\pi$, and $e^{\pi\sqrt{163}}$ are
algebraically independent.  In particular, $e + \pi$ is transcendental (it
generates a field meeting the Nesterenko conditions).  This is an
\emph{unconditional} result, not requiring SC.  Whether it connects directly
to the status of $\tan(1)$ is unclear: Nesterenko's result concerns
$e^{\pi} = (e^{i})^{-i}$ and uses modular functions, whereas $\tan(1)$
involves $e^i$ with $\pi$ absent.  We record it here as a remark, not a
claimed connection.
\end{remark}

\begin{remark}[What SC does not give]
Schanuel's conjecture is not proved.  Proposition~\ref{prop:SCstrength} is
therefore a \emph{conditional strengthening}: it says ``if SC then the
obstruction to $\tan(1) \in \EML{1}$ is deeper than transcendence over
$\mathbb{Q}$.''  The unconditional proof of $\tan(1) \notin \EML{1}$
(Theorem~\ref{thm:DOOR4}) does not require SC.
\end{remark}

%% ------------------------------------------------------------------
\section{Impact on the Depth Hierarchy}
%% ------------------------------------------------------------------

We now record, still conditionally, what Hypothesis~\ref{hyp:H} would imply
for the structure of the EML Atlas depth table.

\begin{conditional}[Depth table impact]\label{cond:E}
\emph{Assuming Hypothesis~\ref{hyp:H}:}
\begin{enumerate}
  \item[(E1)] \textbf{arctan drops from depth 3 to depth 2.}
    Since $\arctan(1) = \pi/4$ and $\pi$ would be EML-constructible
    (Conditional~\ref{cond:C}), the depth-3 tree for $\arctan$ (which uses
    the complex logarithm form $\arctan(x) = \tfrac{1}{2i}\ln\!\left(\tfrac{1+ix}{1-ix}\right)$)
    would be superseded by a depth-2 construction that passes through the
    now-accessible value $\pi/4 = \arctan(1)$ as an anchor.

  \item[(E2)] \textbf{The five-way equivalence collapses.}
    The capstone paper establishes a five-way equivalence among:
    completeness, depth stability, i-unconstructibility (T18), depth-3 ceiling
    for standard functions, and Lindemann–Weierstrass applicability.
    Hypothesis~\ref{hyp:H} breaks T18 (by Conditional~\ref{cond:A}), which by
    the equivalence breaks all five simultaneously.

  \item[(E3)] \textbf{The depth-3 ceiling for standard functions breaks.}
    The current EML Atlas assigns depth $\leq 3$ to all standard elementary
    functions.  Under Hypothesis~\ref{hyp:H}, $\arctan$ would be depth 2, and
    it is unclear whether any function would remain at depth 3 (some might
    drop, breaking the uniformity of the ceiling).

  \item[(E4)] \textbf{The Depth Stability Theorem becomes false.}
    As stated in Conditional~\ref{cond:D}: $\depthC(f) < \depthR(f)$ for
    at least one Atlas function, contradicting DST.
\end{enumerate}
\end{conditional}

\begin{remark}[Why this is not a paradox]
The conditionals above are logically consistent statements about what would be
true \emph{if} H held.  They do not create a paradox because H is false
(Theorem~\ref{thm:DOOR4}).  Their value is diagnostic: they show that
$\tan(1) \in \EML{1}$ is not merely a local curiosity but would require the
entire depth-hierarchy framework to be rebuilt from scratch.  The fact that H
leads to five simultaneous collapses is strong evidence (of a heuristic,
not formal, kind) for why no depth-1 tree can produce $\tan(1)$.
\end{remark}

%% ------------------------------------------------------------------
\section{Summary Table}
%% ------------------------------------------------------------------

\begin{center}
\begin{tabular}{p{5cm}cc}
\hline
\textbf{Statement} & \textbf{Status} & \textbf{Source} \\
\hline
$i \notin \EML{1}$ & Theorem (T18) & Lean-verified \\
Depth Stability Theorem (DST) & Theorem & T18 equivalent \\
$\tan(1)$ is transcendental over $\mathbb{Q}$ & Theorem & Lindemann--Weierstrass \\
$\tan(1) \notin \EML{1}$ & Theorem (DOOR-4) & T18 + contrapositive \\
\hline
$\tan(1) \in \EML{1} \Rightarrow i \in \EML{1}$ & Conditional~A & hypothetical \\
$\tan(1) \in \EML{1} \Rightarrow \sin(1),\cos(1) \in \EML{O(1)}$ & Conditional~B & hypothetical \\
$\tan(1) \in \EML{1} \Rightarrow \pi \in \EML{O(1)}$ & Conditional~C & hypothetical + arcsin depth \\
$\tan(1) \in \EML{1} \Rightarrow$ DST fails & Conditional~D & hypothetical \\
$\tan(1) \in \EML{1} \Rightarrow \arctan$ drops to depth 2 & Conditional~E1 & hypothetical \\
\hline
SC $\Rightarrow \trdeg_{\mathbb{Q}}(e,\cos(1),\sin(1)) \geq 2$ & Conditional (SC) & Schanuel \\
SC $\Rightarrow \tan(1) \notin \overline{\mathbb{Q}(e)}$ & Conditional (SC) & Schanuel + Prop.~5.2 \\
\hline
\end{tabular}
\end{center}

%% ------------------------------------------------------------------
\section*{Notes on Rigour}
%% ------------------------------------------------------------------

The conditional steps (A)–(D) are informal logical derivations, not formal
proofs.  Step~A in particular relies on a depth-counting argument over
$\mathbb{C}$ that has not been formalised in Lean.  The conclusion that
$i \in \EML{1}$ from $\tan(1) \in \EML{1}$ should be understood as: ``under
reasonable assumptions about how algebraic operations over EML-reachable
values compose in depth, $i$ would be EML-reachable.''  The contrapositive
(Theorem~\ref{thm:DOOR4}) is rigorous because it rests on T18 (Lean-verified)
and elementary logic; the precise depth path from $\tan(1)$ to $i$ is not
needed for the validity of the contrapositive, only the existence of \emph{some}
path.

\medskip
\noindent
Formalising Step~A in Lean would require: (1) a definition of complex EML
depth, (2) the algebraic closure lemma (EML-reachable values are closed under
field operations in bounded depth), and (3) the complex-to-real depth transfer
for the specific case of $i$.  These are tractable but unfinished.

\end{document}
